\section*{Chapitre \ref{ch:b2:neurons-synapses} --- Neurones, potentiels d'action et synapses}

\begin{solution}{exo:b2:neurons-synapses:1}
Les \omterm{def:b2:neurons-synapses:neuron}{dendrites} reçoivent les \omterm{def:b2:neurons-synapses:synapse}{synapses}~; le soma contient le noyau et
intègre~; le \omterm{prop:b2:neurons-synapses:integration}{cône d'émergence} décide~; l'\omterm{def:b2:neurons-synapses:neuron}{axone} conduit~; les
terminaisons libèrent le \omterm{def:b2:neurons-synapses:synapse}{neurotransmetteur}. La \omterm{def:b2:neurons-synapses:neuron}{myéline} est la membrane
enroulée d'une cellule gliale, presque dépourvue de cytoplasme~; elle
isole l'\omterm{def:b2:neurons-synapses:neuron}{axone}, si bien que l'influx saute d'un nœud à l'autre et se
propage cent fois plus vite, pour un coût moindre.
\end{solution}

\begin{solution}{exo:b2:neurons-synapses:2}
Dépolarisation jusqu'au seuil~; phase ascendante --- ouverture des
canaux sodiques voltage-dépendants, rétroaction positive~; sommet près
de $E_{\mathrm{Na}}$ --- inactivation des canaux sodiques~; phase
descendante --- ouverture des canaux potassiques voltage-dépendants~;
hyperpolarisation --- canaux potassiques encore ouverts, près de
$E_{\mathrm K}$~; retour au repos à mesure qu'ils se ferment. En
tout-ou-rien parce que la rétroaction positive va jusqu'à son terme ou
ne démarre pas du tout~; à sens unique parce que la membrane qui vient
d'être parcourue est réfractaire (canaux sodiques inactivés).
\end{solution}

\begin{solution}{exo:b2:neurons-synapses:3}
L'influx atteint la terminaison~; des canaux calciques voltage-dépendants
s'ouvrent~; le calcium déclenche la fusion des vésicules~; le
\omterm{def:b2:neurons-synapses:synapse}{neurotransmetteur} diffuse à travers la fente~; il se fixe sur les
récepteurs postsynaptiques~; des canaux s'ouvrent (ou une \omterm{prop:b2:cell-signalling:gpcr}{protéine G}
commute)~; \omterm{def:b2:neurons-synapses:synapse}{potentiel postsynaptique}~; élimination du \omterm{def:b2:neurons-synapses:synapse}{neurotransmetteur}
par une enzyme ou un transporteur~; récupération de la membrane des
vésicules.
\end{solution}

\begin{solution}{exo:b2:neurons-synapses:4}
Électriques~: pas de délai, en général dans les deux sens, pas
d'amplification, pas d'inhibition (le courant ne fait que passer).
Chimiques~: une demi-milliseconde, à sens unique, amplification (un
influx libère des centaines de quanta), inhibition possible (canaux
chlorure ou potassiques), et modifiables.
\end{solution}

\begin{solution}{exo:b2:neurons-synapses:5}
D'après $E = (61.5/z)\log_{10}(c_{\text{ext}}/c_{\text{int}})$~:
potassium \qty{-89}{mV}~; sodium \qty{+61}{mV}~; chlorure, avec $z =
-1$, $-61.5\log_{10} 11 = \qty{-64}{mV}$~; calcium, avec $z = 2$,
$30.75\times 4.3 = \qty{+132}{mV}$. À \qty{-70}{mV}~: le potassium
sort, le sodium entre, le chlorure sort légèrement (la membrane est en
dessous de son \omterm{thm:b2:neurons-synapses:nernst}{potentiel d'équilibre}), le calcium entre fortement.
\end{solution}

\begin{solution}{exo:b2:neurons-synapses:6}
$20 : 1$~: $(20\times(-89) + 61)/21 = \qty{-82}{mV}$, l'état de repos,
près de $E_{\mathrm K}$. $1 : 20$~: $(-89 + 20\times 61)/21 =
\qty{+54}{mV}$, le sommet, près de $E_{\mathrm{Na}}$. $50 : 1$~: $(50\times
(-89) + 61)/51 = \qty{-86}{mV}$, l'hyperpolarisation, avec des canaux
potassiques supplémentaires ouverts.
\end{solution}

\begin{solution}{exo:b2:neurons-synapses:7}
$\sqrt{500} = 22$~: \qty{22}{m/s}. Pour atteindre \qty{100}{m/s}, un
\omterm{def:b2:neurons-synapses:neuron}{axone} amyélinisé aurait besoin de $d = 10^{4}$ \unit{\micro m}, un
centimètre. Un mètre demande \qty{10}{ms} dans la fibre myélinisée,
\qty{45}{ms} dans l'\omterm{def:b2:neurons-synapses:neuron}{axone} de calmar, une seconde entière dans une fibre
de \qty{1}{\micro m}.
\end{solution}

\begin{solution}{exo:b2:neurons-synapses:8}
$m = \ln(200/74) = 1.0$~; $P(1) = e^{-1} = 0.37$, $P(2) = 0.18$, $P(3)
= 0.06$. Réponse moyenne $1.0\times 0.5 = \qty{0.5}{mV}$.
\end{solution}

\begin{solution}{exo:b2:neurons-synapses:9}
$1/\qty{2}{ms} = 500$ influx par seconde. L'intensité du toucher est
codée par la fréquence des influx (et par le nombre de \omterm{def:b2:neurons-synapses:neuron}{neurones}
recrutés)~; les influx eux-mêmes sont identiques.
\end{solution}

\begin{solution}{exo:b2:neurons-synapses:10}
$\lambda = \sqrt{R_m d/4R_i}$~: \qty{0.5}{mm}, \qty{1.6}{mm},
\qty{11}{mm}. Avec la \omterm{def:b2:neurons-synapses:neuron}{myéline}, pour \qty{10}{\micro m}~: $1.6\times\sqrt{200} =
\qty{22}{mm}$. Un internœud de \qty{1}{mm} ne perd que $1 - e^{-1/22}
= \qty{4}{\%}$ du signal, si bien que le nœud suivant est facilement amené
au seuil. À travers \qty{3}{mm} démyélinisés, le signal tombe à
$e^{-3/1.6} = \qty{15}{\%}$ --- pour un influx de \qty{100}{mV}, à peu
près le seuil de \qty{15}{mV}~: la conduction échoue ou devient
incertaine.
\end{solution}

\begin{solution}{exo:b2:neurons-synapses:11}
Tétrodotoxine~: pas de courant sodique, pas d'influx, pas de
transmission. Tétraéthylammonium~: pas de courant potassique, influx
prolongé, davantage de \omterm{def:b2:neurons-synapses:synapse}{neurotransmetteur} libéré par influx. Sans
calcium~: influx normaux, pas de libération, pas de transmission.
Curare~: libération normale, récepteurs bloqués, pas de potentiel de
plaque motrice --- paralysie. Inhibiteur de la cholinestérase~:
l'acétylcholine persiste, les récepteurs restent ouverts, le muscle se
dépolarise puis ne peut plus se repolariser --- paralysie par excès.
Toxine botulique~: les vésicules ne peuvent pas fusionner, pas de
libération --- paralysie.
\end{solution}

\begin{solution}{exo:b2:neurons-synapses:12}
Le délai achète~: la transmission dans un seul sens~; l'amplification
(un seul influx libère des centaines de quanta et peut commander une
cellule plus grande)~; l'inhibition, dont une jonction électrique est
incapable~; et la plasticité, puisque le nombre de quanta et de
récepteurs peut changer avec l'usage, ce qui permet aux circuits
d'apprendre. Les \omterm{def:b2:neurons-synapses:synapse}{synapses} électriques sont utilisées là où la synchronie
et la vitesse comptent davantage que le calcul~: circuits de fuite,
muscle cardiaque, certains réseaux inhibiteurs.
\end{solution}

\begin{solution}{pb:b2:neurons-synapses:1}
\textbf{1.} $E_{\mathrm K} = 61.5\log_{10}(4/140) = \qty{-95}{mV}$~;
$E_{\mathrm{Na}} = 61.5\log_{10}(145/12) = \qty{+67}{mV}$.
\textbf{2.} $(25\times(-95) + 67)/26 = \qty{-89}{mV}$.
\textbf{3.} $E_{\mathrm K} = 61.5\log_{10}(8/140) = \qty{-76}{mV}$~;
$V = (25\times(-76) + 67)/26 = \qty{-71}{mV}$~: plus près du seuil, donc
d'abord plus excitable --- et, si cela persiste, moins, car les canaux
sodiques s'inactivent au niveau dépolarisé.
\textbf{4.} Les gradients s'épuisent en quelques heures, le sodium
entrant et le potassium sortant par les fuites~; le potentiel dérive
vers zéro~; la pompe arrêtée, les anions imperméants de la cellule
attirent le chlorure et l'eau, et elle gonfle.
\textbf{5.} Surface par millimètre $\pi\times 15\times 10^{-6}\times
10^{-3} = \qty{4.7e-8}{m^2}$~; $C = \qty{4.7e-10}{F}$~; $Q = CV =
\qty{4.7e-11}{C}$~: $2.9\times 10^{8}$ ions.
\textbf{6.} Volume par millimètre $\pi(7.5\times 10^{-6})^{2}\times
10^{-3} = \qty{1.8e-13}{m^3} = \qty{1.8e-10}{L}$, contenant $12\times
10^{-3}\times 1.8\times 10^{-10} = 2.1\times 10^{-12}$ mol, soit $1.3\times
10^{12}$ ions~: une variation de $2\times 10^{-4}$.
\textbf{7.} $2.9\times 10^{8}/3 \approx 10^{8}$ ATP par millimètre et par
influx.
\textbf{8.} $0.9/90 = \qty{10}{ms}$ dans chaque sens.
\textbf{9.} Nu~: $\sqrt{1\times 15\times 10^{-6}/4} = \qty{1.9}{mm}$~;
myélinisé~: $\times\sqrt{250} = \qty{31}{mm}$.
\textbf{10.} Nu~: $e^{-1.5/1.9} = 0.45$, \qty{45}{mV}~; myélinisé~:
$e^{-1.5/31} = 0.95$, \qty{95}{mV}. Les deux dépassent le seuil de
\qty{15}{mV}~; mais l'\omterm{def:b2:neurons-synapses:neuron}{axone} nu doit charger toute sa membrane en chemin,
ce qui le rend lent, alors que sous la \omterm{def:b2:neurons-synapses:neuron}{myéline} presque rien n'est perdu
ni chargé.
\textbf{11.} $\sqrt{15} = \qty{3.9}{m/s}$~: \qty{0.23}{s} dans chaque sens,
près d'une demi-seconde pour l'aller-retour --- trop lent pour un réflexe
qui doit rattraper un faux pas.
\textbf{12.} $e^{-4/1.9} = 0.12$~: \qty{12}{mV} au nœud distant, sous le
seuil --- l'influx est bloqué et le réflexe est perdu.
\textbf{13.} Derrière l'influx, les canaux sodiques restent inactivés
pendant une ou deux millisecondes~: la membrane qui vient d'être
parcourue ne peut pas être réexcitée, et l'onde ne peut qu'avancer.
\textbf{14.} $m = \ln(300/111) = 1.0$.
\textbf{15.} $P(1) = 0.37$, $P(2) = 0.18$, $P(3) = 0.06$~: environ 110, 55
et 18 des 300 essais.
\textbf{16.} $1.0\times 0.4 = \qty{0.4}{mV}$.
\textbf{17.} $P(0) = e^{-60} \approx 10^{-26}$~: jamais~; réponse moyenne
\qty{24}{mV}, au-dessus du seuil de \qty{15}{mV}~: une seule \omterm{def:b2:neurons-synapses:synapse}{synapse} de
ce type fait décharger le motoneurone.
\textbf{18.} Vingt \omterm{def:b2:neurons-synapses:synapse}{synapses} actives ensemble additionnent leurs courants
(sommation spatiale), de façon sous-linéaire à mesure que le potentiel
approche du potentiel d'inversion des canaux~; même avec $m = 1$ chacune,
elles donneraient \qty{8}{mV}, et avec un $m$ normal, bien plus que le seuil.
\textbf{19.} Ouvrir des canaux chlorure à $E_{\mathrm{Cl}} = V_{\text{repos}}$
ajoute de la conductance sans déplacer le potentiel~; le courant
excitateur traverse désormais une membrane plus perméable et produit une
dépolarisation plus petite --- une inhibition par court-circuit.
\textbf{20.} $10 + 0.6 + 10 + 0.6 + 5 = \qty{26}{ms}$, contre
\qty{30}{ms} mesurées.
\textbf{21.} L'activation du récepteur lui-même dans le fuseau
neuromusculaire, la propagation du \omterm{prop:b2:neurons-synapses:ap}{potentiel d'action} de la \omterm{def:b2:cell-differentiation:myogenesis}{fibre
musculaire} le long de la fibre (quelques mètres par seconde sur quelques
centimètres), et la latence mécanique avant l'apparition de la force.
\textbf{22.} Chaque \omterm{def:b2:neurons-synapses:synapse}{synapse} ajoute une demi-milliseconde et un
interneurone ajoute une cellule~; mais les interneurones permettent
d'inverser le signal (en inhibant l'antagoniste), de le combiner avec
d'autres et de le faire moduler par le cerveau --- un réflexe que l'on
peut façonner.
\textbf{23.} $90 - 6.2 = \qty{84}{ms}$ de conduction pour \qty{1.8}{m}~:
environ \qty{21}{m/s}.
\textbf{24.} La moitié du courant entrant à chaque tension~: le seuil
s'élève, l'influx est plus petit et plus lent, le facteur de sécurité à
chaque nœud diminue, et l'influx peut échouer~; le réflexe s'affaiblit
ou disparaît.
\textbf{25.} \omterm{thm:b2:neurons-synapses:chord}{Potentiel de repos} \qty{-89}{mV}~; \qty{10}{ms} dans chaque
sens~; $m = 1.0$~; latence calculée \qty{26}{ms}.
\end{solution}
